%&bericht

%%%%%%%%%%%%%%%%%%%%%%%%%%%%%%%%%%%%%%%%%%%%%%%%%%%%%%%%%%%%%%%%%%%%%%%%%%%%%%%
%% Descr:       Vorlage für Berichte der DHBW-Karlsruhe
%% Author:      Prof. Dr. Jürgen Vollmer, juergen.vollmer@dhbw-karlsruhe.de
%% $Id: bericht.tex,v 1.30 2023/07/25 11:03:01 vollmer Exp $
%%  -*- coding: utf-8 -*-
%%%%%%%%%%%%%%%%%%%%%%%%%%%%%%%%%%%%%%%%%%%%%%%%%%%%%%%%%%%%%%%%%%%%%%%%%%%%%%%

\documentclass[
   ngerman          % neue deutsche Rechtschreibung
  ,a4paper          % Papiergrösse
% ,twoside          % Zweiseitiger Druck (rechts/links)
% ,10pt             % Schriftgrösse
  ,11pt
% ,12pt
  ,pdftex
%  ,disable         % Todo-Markierungen auschalten
  ,dvipsnames,table
]{report}

% Bitte die Codierung Ihrer Dateien auswählen:
% \usepackage[latin1]{inputenc}    % Für UNIX mit ISO-LATIN-codierten Dateien
% \usepackage[applemac]{inputenc}  % Für Apple Mac
% \usepackage[ansinew]{inputenc}   % Für Microsoft Windows
\usepackage[utf8]{inputenc}        % UTF-8 codierte Dateien
                                   % Dieses Dokument ist unter Unix erstellt, daher
                                   % wird diese Input-Codierung benutzt.
\usepackage[acronym]{glossaries}
\usepackage{bericht}
\usepackage{pdfpages}
\makeglossaries
\loadglsentries{glossaries.tex}

%% ACHTUNG, wenn man eine eigene Formatdatei (bericht.fmt) benutzt, werden Änderungen an bericht.sty
%% erst wirksam, wenn die Format-Datei neu erzeugt wurde!!!
%% Genauer alle Änderungen, die textuell vor der nächsten Zeile ".... endofdump...." stehen
%% werden erst wirksam, wenn die Formatdatei neu erzeugt wurde
\csname endofdump\endcsname

%%%%%%%%%%%%%%%%%%%%%%%%%%%%%%%%%%%%%%%%%%%%%%%%%%%%%%%%%%%%%%%%%%%%%%%%%%%%%%%
%% Angaben zur Arbeit
%%%%%%%%%%%%%%%%%%%%%%%%%%%%%%%%%%%%%%%%%%%%%%%%%%%%%%%%%%%%%%%%%%%%%%%%%%%%%%%

\newcommand{\Autor}{Benjamin Dangl}
\newcommand{\MatrikelNummer}{99999}
\newcommand{\Kursbezeichnung}{tinf21b4}

\newcommand{\FirmenName}{SICK AG}
\newcommand{\FirmenStadt}{Waldkirch}

% Falls es kein Firmenlogo gibt:
\newcommand{\FirmenLogoDeckblatt}{}

\newcommand{\BetreuerDHBW}{Mirko Dostmann}

%%%%%%%%%%%%%%%%%%%%%%%%%%%%%%%%%%%%%%%%%%%%%%%%%%%%%%%%%%%%%%%%%%%%%%%%%%%%%%%%%%%%%

% \newcommand{\Abschluss}{Bachelor of Engineering}
\newcommand{\Abschluss}{Bachelor of Science}

% \newcommand{\Studiengang}{Informatik / Informationstechnik}
\newcommand{\Studiengang}{Informatik}

\hypersetup{%%
  pdfauthor={\Autor},
  pdftitle={Technische Dokumentation},
  pdfsubject={Volunteer Organizer},
}

%%%%%%%%%%%%%%%%%%%%%%%%%%%%%%%%%%%%%%%%%%%%%%%%%%%%%%%%%%%%%%%%%%%%%%%%%%%%%%%

% Wenn \includeonly{..} benutzt wird, werden nur diese Kaptitel ausgegeben.
\includeonly{
  glossaries
 ,kapitel1
 ,kapitel2
 ,kapitel3
 ,kapitel4
 ,kapitel5
 ,kapitel6
 ,appendix
}

%%%%%%%%%%%%%%%%%%%%%%%%%%%%%%%%%%%%%%%%%%%%%%%%%%%%%%%%%%%%%%%%%%%%%%%%%%%%%%%

% Benutzt man das "biblatex"-Paket, dann muß das hier stehen:
% siehe auch die mit BIBLATEX markierten Zeilen in bericht.sty
\bibliography{bericht}

\begin{document}

\pagenumbering{roman}
% Seitennummern in römischen bzw. arabischen Ziffern:
% https://tex.stackexchange.com/questions/226864/start-toc-from-chapter-1-and-also-page-numbering
% style can be any of these:
%
%    arabic: arabic numerals
%    roman: lowercase roman numerals
%    Roman: uppercase roman numerals
%    alph: lowercase letters
%    Alph: uppercase letters

%%%%%%%%%%%%%%%%%%%%%%%%%%%%%%%%%%%%%%%%%%%%%%%%%%%%%%%%%%%%%%%%%%%%%%%%%%%%%%%

\begin{titlepage}
  \begin{center}
    \vspace*{-2cm}
    \FirmenLogoDeckblatt\hfill\includegraphics[width=4cm]{assets/figures/dhbw-logo.png}\\[2cm]
    {\Huge Technische Dokumentation}\\[1cm]
    {\Huge\scshape Volunteer Organizer}\\[1cm]
    {\large für die Prüfungsleistung}\\[0.5cm]
    {\Large Advanced Software Engineering}\\[0.5cm]
    {\large des Studienganges \Studiengang}\\[0.5cm]
    {\large an der}\\[0.5cm]
    {\large Dualen Hochschule Baden-Württemberg Karlsruhe}\\[0.5cm]
    {\large von}\\[0.5cm]
    {\large\bfseries \Autor}
    \vfill
  \end{center}
  \begin{tabular}{l@{\hspace{2cm}}l}
    Matrikelnummer                & \MatrikelNummer  \\
    Kurs                          & \Kursbezeichnung \\
    Gutachter der Studienakademie & \BetreuerDHBW    \\
  \end{tabular}
\end{titlepage}

%%%%%%%%%%%%%%%%%%%%%%%%%%%%%%%%%%%%%%%%%%%%%%%%%%%%%%%%%%%%%%%%%%%%%%%%%%%%%%%

%%%%%%%%%%%%%%%%%%%%%%%%%%%%%%%%%%%%%%%%%%%%%%%%%%%%%%%%%%%%%%%%%%%%%%%%%%%%%%%
%% Descr:       Vorlage für Berichte der DHBW-Karlsruhe, Erklärung
%% Author:      Prof. Dr. Jürgen Vollmer, vollmer@dhbw-karlsruhe.de
%% $Id: erklaerung.tex,v 1.11 2020/03/13 14:24:42 vollmer Exp $
%% -*- coding: utf-8 -*-
%%%%%%%%%%%%%%%%%%%%%%%%%%%%%%%%%%%%%%%%%%%%%%%%%%%%%%%%%%%%%%%%%%%%%%%%%%%%%%%

% In Bachelorarbeiten muss eine schriftliche Erklärung abgegeben werden.
% Hierin bestätigen die Studierenden, dass die Bachelorarbeit, etc.
% selbständig verfasst und sämtliche Quellen und Hilfsmittel angegeben sind. Diese Erklärung
% bildet das zweite Blatt der Arbeit. Der Text dieser Erklärung muss auf einer separaten Seite
% wie unten angegeben lauten.

\newpage
\thispagestyle{empty}
\begin{framed}
  \begin{center}
    \Large\bfseries Erklärung
  \end{center}
  \medskip
  \noindent
  % siehe §5(3) der \enquote{Studien- und Prüfungsordnung DHBW Technik} vom 29.\,9.\,2017 und Anhang 1.1.13
  Ich versicheree hiermit, dass ich meine Technische Dokumentation der Software \enquote{Volunteer Organizer}
  selbstständig verfasst und keine anderen als die angegebenen Quellen und Hilfsmittel benutzt habe. % Ich versichere zudem, dass die eingereichte elektronische Fassung mit der gedruckten Fassung übereinstimmt.
  \vspace{3cm}\\
  \noindent
  \underline{\hspace{4cm}}\hfill\underline{\hspace{6cm}}\\
  Ort~~~~~Datum\hfill Unterschrift\hspace{4cm}
\end{framed}

% \vfill
% \emph{Sofern  vom Dualen Partner ein Sperrvermerk gewünscht wird, ist folgende Formulierung
% zu verwenden:}
% \begin{framed}
% \begin{center}
% \Large\bfseries Sperrvermerk
% \end{center}
% \medskip
% \noindent
% Der Inhalt dieser Arbeit darf weder als Ganzes noch in Auszügen Personen
% außerhalb des Prüfungsprozesses und des Evaluationsverfahrens zugänglich gemacht
% werden, sofern keine anderslautende Genehmigung vom Dualen Partner vorliegt.
% \end{framed}

%%%%%%%%%%%%%%%%%%%%%%%%%%%%%%%%%%%%%%%%%%%%%%%%%%%%%%%%%%%%%%%%%%%%%%%%%%%%%%%
\endinput
%%%%%%%%%%%%%%%%%%%%%%%%%%%%%%%%%%%%%%%%%%%%%%%%%%%%%%%%%%%%%%%%%%%%%%%%%%%%%%%


%%%%%%%%%%%%%%%%%%%%%%%%%%%%%%%%%%%%%%%%%%%%%%%%%%%%%%%%%%%%%%%%%%%%%%%%%%%%%%%

\begin{abstract}

\end{abstract}

\newcounter{savepageRoman} % needed to continue after text

\newpage
\tableofcontents           % Inhaltsverzeichnis hier ausgeben
\listoffigures             % Liste der Abbildungen
\listoftables              % Liste der Tabellen
\lstlistoflistings         % Liste der Listings
\listofequations           % Liste der Formeln

% Jetzt kommt der "eigentliche" Text
\newacronym{KI}{KI}{Künstliche Intelligenz}
\newacronym{RLHF}{RLHF}{Reinforcement Learning with Human Feedback}              % Abkürzungsverzeichnis
\setcounter{savepageRoman}{\value{page}}
\pagenumbering{arabic}
% https://tex.stackexchange.com/questions/226864/start-toc-from-chapter-1-and-also-page-numbering

\chapter{Domain Driven Design}
\label{chapter:domain_driven_design}

\section{Analyse der Ubiquitous Language}
\label{section:analyse_der_ubiquitous_language}

Die Ubiquitous Language dieses Projekts richtet sich nach der üblichen Sprache in der Vereinsdomäne. Die folgenden Begriffe sind speziell aus der Subdomäne \glqq Veranstaltungen\grqq{} entnommen und werden im weiteren Verlauf des Projekts verwendet.

\paragraph{\glqq Veranstaltung\grqq{} (Event)}
Dieser begriff bezeichnet eine Vereinsveranstaltung, die Freiwillige für die Durchführung benötigt. Eine Veranstaltung findet in einem bestimmten Zeitraum statt und hat einen Veranstaltungsort. Außerdem gibt es verschiedene Aufgaben, die in einem bestimmten Zeitfenster durch Freiwillige erledigt werden müssen. Durch die verschiedenen Aufgabenprofile sind auch verschiedene Eigenschaften von Freiwilligen gefragt. Eine Veranstaltung wird durch mindestens einen Organisator organisiert.

\paragraph{\glqq Freiwilliger\grqq{} (Volunteer)}
Dieser Begriff bezeichnet eine Person, die bei einer Veranstaltung des Vereins hilft. Ein Freiwilliger kann verschiedene Eigenschaften besitzen, die für die Durchführung einer Veranstaltung benötigt werden. Ein Freiwilliger kann auch Mitglied im Verein sein.

\paragraph{\glqq Aufgabe\grqq{} (Task)}
Dieser Begriff bezeichnet eine konkrete Tätigkeit, die bei einer Veranstaltung in einem bestimmten Zeitraum durchgeführt werden muss. Außerdem kann es sein, dass eine Aufgabe bestimmte Anforderungen an die Eigenschaften von Freiwilligen hat. Eine Aufgabe wird durch mindestens einen Freiwilligen erledigt.

\paragraph{\glqq Eigenschaft\grqq{} (Feature)}
Dieser Begriff bezeichnet eine Fähigkeit, die ein Freiwilliger besitzt. Neben technischen oder fachlichen Fähigkeiten kann dies auch der Besitz eines Führerscheins oder der Besitz eines Schlüssels für eine bestimmte Räumlichkeit sein.

\paragraph{\glqq Mitglied\grqq{} (Member)}
Dieser Begriff bezeichnet eine Person, die Mitglied im Verein ist. Prinzipiell stehen Mitglieder als potentielle Freiwillige zur Verfügung. Sie haben also auch verschiedene Eigenschaften.

\paragraph{\glqq Veranstaltungsort\grqq{} (Location)}
Dieser Begriff bezeichnet den Ort, an dem eine Veranstaltung stattfindet. Der Ort kann auch \glqq virtuell\grqq{} sein, wenn die Veranstaltung online stattfindet.

\paragraph{\glqq Organisator\grqq{} (Organizer)}
Dieser Begriff bezeichnet ein Vereinsmitglied, das eine Veranstaltung organisiert. Zur Organisation gehört das Festlegen von Ort und Zeit, das Festlegen von Aufgaben und das Einladen von Freiwilligen.

\section{Analyse und Begründung der verwendeten Muster}
\label{section:analyse_und_begrundung_der_verwendeten_muster}


\chapter{Programming Principles}
\label{chapter:programming_principles}

\section{Analyse und Begründung der ausgewählten Prinzipien}

\todo{Analyse und Begründung für 5
  (fünf) der vorgestellten Prinzipien aus
  • SOLID
  • GRASP
  • DRY}
\chapter{Unit Tests}
\label{chapter:unit_tests}

\todo{Fill content}
\chapter{Refactoring}
\label{chapter:refactoring}

\section{Code Smells}
\label{section:code_smells}

\todo{Code Smells finden und
  Refactoring begründen}
\chapter{Entwurfsmuster}
\label{chapter:entwurfsmuster}

\section{Begründung des Einsatzes}

\todo{Einsatz ausführlich begründen}

% Ab hier beginnt der Anhang
\pagenumbering{roman}
\setcounter{page}{\value{savepageRoman}}
\appendix
\addcontentsline{toc}{chapter}{Anhang}
\include{appendix}
\addcontentsline{toc}{chapter}{Index}
\printindex
\printglossary[type=\acronymtype]

\addcontentsline{toc}{chapter}{Literaturverzeichnis}

% Haben Sie das "biblatex"-Paket nicht installiert, benutzen Sie folgendes:
% Ohne das "biblatex"-Paket (s. bericht.sty) produziert folgendes
% "deutsche" Zitate in Literaturverzeichnissen gemaß der Norm DIN 1505,
% Teil 2 vom Jan. 1984.
% Die Zitatmarken werden alphabetisch nach Verfassern
% sortiert und sind durch abgekürzte Verfasserbuchstaben plus
% Erscheinungsjahr in eckigen Klammern gekennzeichnet.

% \bibliographystyle{alphadin}
% \bibliography{bericht}

%%%%%%%%%%%%%%%%%%%%%%%%%%%%%%%%%%%%%%%5
% BIBLATEX
% Benutzt man das "biblatex"-Paket, muß man folgendes schreiben:
\def\refname{Literaturverzeichnis}
\printbibliography
%%%%%%%%%%%%%%%%%%%%%%%%%%%%%%%%%%%%%%%

\newpage
\addcontentsline{toc}{chapter}{Liste der ToDo's}
\listoftodos[Liste der ToDo's]


\end{document}
